\chapter{时间序列与统计套利分层答案}
\begin{enumerate}[leftmargin=*]
  \item 弱白噪声只要求零均值、常方差和跨期不相关；创新是相对当前信息集不可预测的误差；错设残差可能保留显著自相关，因此不能把三者混用。
  \item 高相关只描述样本共同变化；协整要求非平稳序列存在平稳线性组合；可交易还需时点正确、关系稳定、成本后为正且有可执行容量。
  \item 因 $X_t=0.6X_{t-1}$，分子等于 $0.6\sum X_{t-1}^2$，故斜率为 $0.6$。均值模型残差仍遵循同一动态，残差一阶系数超过阈值即拒绝。
  \item $K=1/(1+0.25)=0.8$，$m^+=0.8\times1.2=0.96$，$P^+=(1-0.8)\times1=0.2$。
  \item 应对每个阈值保存检测索引、真实变点、延迟、提前误报和是否漏检；只报最佳阈值属于选择后汇报。
  \item 新交易必须先写独立解析/枚举/第二实现参考账本，再修改 fixture、哈希和 oracle；毛收益与每类成本分别列示。
  \item 随机切分破坏因果时间方向；全样本标准化把未来均值与方差送入过去。二者都使验证集不再模拟当时可获得的信息。
  \item 例如连续三期残差越界或在线状态概率低于阈值即停用，阈值与恢复条件在交易前事先确定；平滑状态只能用于事后诊断。
\end{enumerate}

\section*{独立 oracle 核对表}
AR(1) 系数为 $0.6$，错设残差的一阶相关为 $1$。

MA 预测为 $0.4$；VAR 两维预测为 $(0.7,0.8)$；GARCH 方差为 $0.5$。Kalman 增益、后验均值和后验方差为 $0.8,0.96,0.2$。

已知变点索引为 4，检测延迟和误报均为 0；低阈值有 2 个提前误报，高阈值漏检。协整价差一阶相关约为 $-0.597614$，漂移伪价差被拒绝。毛收益为 $0.05$，成本后净收益为 $0.04$。

\section*{逐题推导与核对}
\begin{enumerate}[leftmargin=*]
  \item 白噪声只要求二阶不相关；创新还要求对当前信息集条件均值为零；错设残差可能仍有自相关、条件异方差或非正态。先写信息集和检验对象，才能决定哪个假设被检验。
  \item 高相关只描述同步变化；协整要求非平稳水平的某个组合平稳；可交易价差还需时点正确、关系稳定、成本后有边际、容量和停用规则。把三个词混用会把统计关系误写成交易结论。
  \item 若 $X_t=0.6X_{t-1}$，OLS 分子为 $\sum X_{t-1}X_t=0.6\sum X_{t-1}^2$，分母为 $\sum X_{t-1}^2$，所以 $\widehat\phi=0.6$。对残差再做一阶回归，错设时相关约为 $1$，说明均值结构未被正确移除。
  \item 预测方差 $P^-=1$，观测方差 $S=P^-+R=1.25$，增益 $K=P^-/S=0.8$；创新为 $1.2$，所以后验均值 $0+0.8(1.2)=0.96$，后验方差 $(1-0.8)1=0.2$。每个数字都来自同一时间点的信息集。
  \item 对每个变点阈值冻结检测规则，记录真实变点索引 4、首个触发索引、提前触发数、延迟和漏检。低阈值的 2 个提前误报与高阈值漏检不能只用平均准确率隐藏；应报告完整阈值表并注明选择前是否固定。
  \item 第三笔交易先由独立参考账本列 target、order、fill、position，再计算毛收益 $0.05$ 和成本后 $0.04$。程序输出逐列比较，若只比较总收益，正负成本错误可能互相抵消。
  \item 随机切分把未来状态混入训练，全样本标准化使用未来均值/方差；修复为按时间拟合 scaler、purge 标签并只在后移窗测试。验证报告需列出切分索引和可用时点。
  \item 例如连续三期残差超过阈值、协整检验失效或在线状态概率低于阈值时停用；阈值、窗口、降仓比例和恢复条件在交易前冻结。平滑状态只能用于事后诊断，不能回填停用时刻。
\end{enumerate}
